\scanpage{27}\mspage{25}
voir encore que $\varphi$ et $\psi$ sont homotopes dans la
catégorie des foncteurs $\unsure{\mathscr{L}|\mathcal{L}}_4\unsure{\mathord{-}|\to}\mathord{\tikz[baseline=0pt,x=1ex,y=1ex]{\draw[line width=.09ex,line join=round] (0,0) -- (.95,1.30) -- (1.52,0) -- cycle; \draw[line width=.09ex] (.32,0) -- (1.02,1.12);}}_2$ \deleted{dont la} qui
sont compatibles avec la source, et le but.

On voit donc que notre lemme d'homotopie
de tantôt n'était pas assez général. Il faudrait
définir les \uline{catégories \unsure{intervalles|intervalle} généralisées}
de longueur $n$ quelc. ($n=\mathrm{nb}$ de flèches bout
à bout), une catégorie d'int. gén. \inserted{\uline{\unsure{standart|standard}}} de long.
$n$ a pour \deleted{\unsure{objets les}} \inserted{ens.} d'objets l'int. d'entiers
\[
[0,\deleted{\unsure{2}}n],
\]
et pour tout couple d'entiers consécutifs
$i-1,i\deleted{\unsure{+1}}$ \deleted{\unsure{qui sont}} \unsure{en} objets ($\unsure{0|1}\leq i\leq n$), il y a
exactement une flèche qui relie $i,i+1$. \deleted{\unsure{Et}}

On affecte l'intervalle $]i,i+1]$ du signe $+$ si
$i\deleted{\leftarrow}\to i+1$ (\inserted{flèche en} sens direct), du signe \unsure{$i$|$-$|$\div$} si
$i\leftarrow i+1$ (flèche en sens rétrograde) -- donc
les catégories \unsure{standart|standard} sont définies, pour $n$ fixé,
par un système de signes $+$ ou $-$, dont on affecte
les $n$ ``\unsure{tronçons}'' consécutifs $\deleted{\unsure{A[i,i+1]}}\inserted{A_i=}[i-1,i]$ ($1\leq i\leq n$)
de l'intervalle d'entiers $[\unsure{0|1},n]$. Il y en a
donc exactement $2^n$.

On s'attend à ce que l'énoncé suivant soit
vrai :

\uline{Lemme} \quad Soient $I,I'$ deux intervalles généralisés,
munis d'objets source et objet but (de sorte que
$I,I'$ sont \uline{canoniquement} isom. aux int. gén.
\unsure{standart|standard} correspondants). Soit $\underline{\mathrm{Hom}}^!(I,I')$
% archive sheet number: 26
